%---------------------------------------------------------------------------%
%->> 封面信息及生成
%---------------------------------------------------------------------------%
%-
%-> 中文封面信息
%-
\confidential{}% 密级：只有涉密论文才填写
\schoollogo{scale=0.2}{ShanghaiTechLogo}% 校徽
%======论文题目 页眉显示 务必准确========
\title{基于物理信息神经网络与大语言模型的统一期权定价框架}% 论文中文题目
\author{朱文龙}% 论文作者
\ID{2022XXXXXX}
\entranceYear{2022}
\advisor{指导教师姓名}% 指导教师：姓名
\advisorsec{}% 第二指导老师：按情况填写
\degree{学士}% 学位：学士、硕士、博士
\degreetype{工学}% 学位类别：理学、工学、工程、医学等
\major{计算机科学与技术}% 二级学科专业名称
\institute{上海科技大学信息科学与技术学院}% 院系名称
\chinesedate{二零二六年六月}% 毕业日期：夏季为6月、冬季为12月
%-
%-> 英文封面信息
%-
\englishtitle{A Unified Option Pricing Framework Based on Physics-Informed Neural Networks and Large Language Models}% 论文英文题目
\englishauthor{Zhu Wenlong}% 论文作者
\englishadvisor{Supervisor: Professor Name}% 指导教师
\englishdegree{Bachelor}% 学位：Bachelor, Master, Doctor。封面格式将根据英文学位名称自动切换，请确保拼写准确无误
\englishdegreetype{Engineering}% 学位类别：Philosophy, Natural Science, Engineering, Economics, Agriculture 等
\englishthesistype{thesis}% 论文类型： thesis, dissertation
\englishmajor{Computer Science and Technology}% 二级学科专业名称
\englishinstitute{School of Information Science and Technology}% 院系名称
\englishdate{June 2026}% 毕业日期：夏季为June、冬季为December
%-
%-> 生成封面
%-
%-======================================================
% 封面和声明页格式不符合教务处要求，因此此处不生成pdf格式的 ||
% 论文封面和声明页，请各位使用word模板中的相关内容。       ||
%=======================================================
%\maketitle% 生成中文封面
%\makeenglishtitle% 生成英文封面
%-
%-> 作者声明
%-
%\makedeclaration% 生成声明页
%-
%-> 中文摘要
%-

\chapter*{摘\quad 要}\chaptermark{摘\quad 要}% 摘要标题
\setcounter{page}{1}% 开始页码
\pagenumbering{Roman}% 页码符号

本文提出并实现了一个基于物理信息神经网络（PINN）与大语言模型（LLM）的统一期权定价框架。针对 Black-Scholes-Merton（BSM）、常弹性方差（CEV）和 Heston 三大经典期权定价模型，分别构建了对应的 PINN 求解器，将 PDE 残差直接嵌入神经网络损失函数，无需网格剖分即可对任意输入即时推断。对于 Heston 模型，采用 ICPINN 硬约束输出变换，确保终值条件在机器精度下精确满足。在网络架构上，BSM 和 CEV 模型使用门控网络（浅层分支与深层分支并联），Heston 模型使用深层全连接网络。在此基础上，引入 DeepSeek 大语言模型作为智能路由层，通过结构化 JSON 输出实现自然语言到期权参数的可靠转换，支持中英文输入。实验结果表明，BSM-PINN 的平均绝对误差（MAE）为 0.056，CEV-PINN 的 MAE 为 1.099，Heston-PINN 的 MAE 为 2.357；LLM 路由层在测试用例上实现了 100\% 的模型选择准确率。本文验证了 PINN 与 LLM 结合用于期权定价的可行性，为智能金融计算提供了一种新的实现路径。

\keywords{物理信息神经网络，期权定价，Black-Scholes 模型，Heston 模型，大语言模型}
%-
%-> 英文摘要
%-
\chapter*{Abstract}\chaptermark{Abstract}% 摘要标题

This thesis proposes and implements a unified option pricing framework based on Physics-Informed Neural Networks (PINN) and Large Language Models (LLM). For three classical option pricing models — Black-Scholes-Merton (BSM), Constant Elasticity of Variance (CEV), and Heston — we construct corresponding PINN solvers that embed PDE residuals directly into the neural network loss function, enabling mesh-free inference at arbitrary inputs. For the Heston model, we adopt the ICPINN hard-constraint output transformation to ensure exact satisfaction of the terminal condition at machine precision. BSM and CEV models use a gated network architecture (parallel shallow and deep branches), while the Heston model uses a deeper fully-connected network. On top of this, we integrate the DeepSeek large language model as an intelligent routing layer, converting natural language descriptions into structured option parameters via JSON-constrained output, supporting both Chinese and English input. Experimental results show that BSM-PINN achieves a mean absolute error (MAE) of 0.056, CEV-PINN achieves MAE of 1.099, and Heston-PINN achieves MAE of 2.357; the LLM routing layer achieves 100\% model selection accuracy on test cases. This work validates the feasibility of combining PINN and LLM for option pricing, providing a new implementation path for intelligent financial computation.

\englishkeywords{Physics-Informed Neural Networks, Option Pricing, Black-Scholes Model, Heston Model, Large Language Models}
%---------------------------------------------------------------------------%
